% abstract.tex - Abstract della tesi
\chapter{Abstract}

L'ottimizzazione del \textit{Length of Stay} (LOS) ospedaliero è una sfida fondamentale per conciliare efficienza economica e qualità delle cure. Questa tesi, sviluppata all'interno del progetto TEXLOS \cite{texlos2026}, propone un sistema innovativo d'intelligenza artificiale per la classificazione retrospettiva e l'individuazione dei colli di bottiglia logistici responsabili di degenze critiche ($\ge 20$ giorni) partendo dai \textit{Log degli Eventi} clinici.

L'approccio metodologico supera i limiti del \textit{Process Mining} tradizionale introducendo il paradigma dello \textbf{Storytelling}: una traduzione semantica dei pattern cronologici in narrazioni in linguaggio naturale. Tali testi alimentano un classificatore Transformer (\texttt{bert-base-uncased}), addestrato mediante \textit{Focal Loss} per mitigare il severo sbilanciamento dei dati ospedalieri.

Avendo accertato l'apprendimento strutturale del modello, viene risolto il problema clinico della scatola nera ("black box"). Si applica un innovativo meccanismo \textsl{adattivo} di \textbf{Integrated Gradients Adattivi} -- fuso alla ricomposizione semantica intelligente dei token frammentati da BERT -- per interrogare a posteriori le scelte della rete. Assegnando inequivocabili "punteggi d'impatto" univoci direttamente alle specifiche azioni nosocomiali storiche, si fornisce ai decisori medici un cruscotto diagnostico di processo (heatmap) matematicamente rigoroso per risolvere empiricamente le inefficienze ospedaliere.
